\subsection{Capacity and measured repair of buffered disks}
\label{sec:collar-capacity}

A logarithmic buffer outside a protected core consumes auxiliary
volume. Fix one source of Theorem~\ref{thm:band-charge}, write
\(\mu_q=dV_q\), and suppose \(Y-6\ell/\sqrt{2}\ge\kappa>0\).
For a compact core \(C_j\) take nested exterior marked annuli
\(A_{j\ell}\), with disjoint interiors and smoothly extending faces.
Assume that the whole region through the outermost face lies in a
single comparison ball of radius at most one in the spectral-median
gauge. Let \(b_{j\ell}>0\) be each annulus's Euclidean logarithmic
width and \(\mathcal L_j=\sum_\ell b_{j\ell}\). In its marking the
determinant-one metric has smallest eigenvalue at least
\(\Lambda^{-1}\). Then
\begin{equation}\label{eq:marked-collar-capacity}
 \sum_j\mathcal L_j^3\mu_q(C_j)
 \le \frac{72\pi^2\Lambda^2}{\kappa^2}
          \sum_{j,\ell}\mu_q(A_{j\ell}).
\end{equation}
If all annuli are disjoint and avoid all the cores, the last sum
is at most \(1-\mu_q(\bigcup_j C_j)\), for pairwise disjoint cores.

Indeed, use one cutoff on each stack, equal to one through the inner
face and zero beyond the outer face. Across a layer let it drop by
\(b_{j\ell}/\mathcal L_j\) linearly in \(\log |x|\), keeping it constant
in the gaps. Its fourth-power Dirichlet integral is conformally
invariant. In determinant-one coordinates its radial derivative
has Euclidean size \((\mathcal L_j |x|)^{-1}\) on the layer, and
the metric costs at most \(\Lambda^2\) in the fourth power.
The summed integral is therefore at most
\(2\pi^2\Lambda^2/\mathcal L_j^3\). The cutoff's nonzero set is
in a comparison ball of at most half the total spectral mass.
The small-support Yamabe inequality in
Theorem~\ref{thm:band-charge} and H\"older give
\[
 \kappa\,\mu_q(C_j)^{1/2}
 \le 6\left(\int |d\chi_j|_q^4\,d\mu_q\right)^{1/2}
          \left(\sum_\ell\mu_q(A_{j\ell})\right)^{1/2}.
\]
Squaring and adding proves \eqref{eq:marked-collar-capacity}.
On a Euclidean annulus the log-radial cutoff attains the minimum
fourth-power energy for a unit drop. H\"older on each ray gives
the same lower bound \(2\pi^2/\mathcal L^3\) after integration over
the unit sphere.

In particular, suppose the protected cores are all the original E58
disks and cover auxiliary volume greater than \(1-d_0\). Disjoint
stacks each containing \(N\) annuli of width two would require
\[
 N^3(1-d_0)<\frac{9\pi^2\Lambda^2}{\kappa^2}\,d_0 .
\]
Thus even one outward annulus of fixed width on every original disk
is excluded at sufficiently small \(d_0\) and fixed \(\Lambda\),
when the cutoffs stay in the stated comparison balls. The estimate
also holds when separate markings supply the distinct layers.
A smaller protected core needs its own measured charge.

\paragraph{A refined allowance in the marked chart.}
On one fixed-width annulus of
Proposition~\ref{prop:disk-collar-comparison}, let
\(d=\|h\|_{2,A}\), retaining \(a=\|h\|_{\infty,A}\),
\(w^2=4(E_++E_-)(g;A)\) and \(\tau=\|DG(m)\|_{2,A}\).
Put \(b=w+\tau+d\). For the same fixed cutoff \(\psi\), now set
\[
 \beta=C(ab+\sqrt{db}+d),\qquad
 n_s=\exp((1-s\psi)h).
\]
For small universal \(a\), in the chiral density representation of
Theorem~\ref{thm:simultaneous} one has
\begin{equation}\label{eq:weighted-collar-allowance}
 \|F_\pm(n_s)-(1-s\psi)F_\pm(m)\|_{2,A}\le s\beta .
\end{equation}
The estimate also holds on either half. Let \(\chi\) be a fixed
interior cutoff equal to one on the support of \(\psi\), and write
\(X=\|\chi^2 D^2h\|_2\), \(Z=\|\chi Dh\|_2\) and
\(J=\|\chi Dh\|_4^2\). The joint elliptic estimate of
Proposition~\ref{prop:disk-collar-comparison} gives
\(X\le C(b+Z+aX+J)\). Integration by parts gives
\(Z^2\le C(dX+dZ)\) and
\(J^2\le C(aXJ+\|h\|_4 J^{3/2})\). For the last term one uses
H\"older on \(h\,D\chi\) and \(|\chi Dh|^3\).
Since \(\|h\|_4^2\le ad\), it follows that
\(Z\le C(\sqrt{dX}+d)\) and \(J\le Ca(X+d)\).
Absorption gives \(X\le Cb\).

Write \(F_\pm(\exp h)=P_\pm(h)D^2h+Q_\pm(h)[Dh,Dh]\) with
smooth coefficients, and set \(\vartheta=1-s\psi\).
The coefficient changes in the density difference are
\(P_\pm(\vartheta h)-P_\pm(h)=O(sa)\) and
\(\vartheta^2Q_\pm(\vartheta h)-\vartheta Q_\pm(h)=O(s)\).
Derivatives landing on \(\vartheta\) cost \(Cs(Z+d)\) in
\(L^2\). This gives \(Cs(aX+J+Z+d)\), proving
\eqref{eq:weighted-collar-allowance}. Restoring the source multiplier
preserves the chiral energies.

\paragraph{Signed path and exterior costs.}
For each selected disjoint patch, let \(B_j\) be the entire raw
disk, \(A_j\) the annulus and \(A_j^{\rm out}\) its exterior half.
For either chirality write
\[
 \begin{gathered}
 S_j^\pm=\lambda^\pm(A_j),\qquad
 S_{j,{\rm out}}^\pm=\lambda^\pm(A_j^{\rm out}),\\
 H_j^\pm=\int_{A_j}(1-\psi)^2d\lambda^\pm,\qquad
 H_{j,{\rm out}}^\pm=\int_{A_j^{\rm out}}(1-\psi)^2d\lambda^\pm .
 \end{gathered}
\]
Using \eqref{eq:weighted-collar-allowance}, put
\begin{align}
 U_j^\pm&=\max\{0,(\sqrt{H_j^\pm}+\beta_j)^2-S_j^\pm\},
       \label{eq:weighted-collar-path}\\
 V_j^\pm&=(\sqrt{H_{j,{\rm out}}^\pm}+\beta_{j,{\rm out}})^2
                                      -S_{j,{\rm out}}^\pm .
       \label{eq:weighted-collar-exterior}
\end{align}
Here \(\beta_{j,{\rm out}}\) may use the restricted error, or the
full \(\beta_j\). The \(L^2\) norm of the leading joined-path density
is the norm of an affine function of \(s\). Adding the allowance
\(s\beta_j\) gives a nonnegative convex function. Its square is
bounded by the larger endpoint square, which proves
\eqref{eq:weighted-collar-path}. The same triangle inequality
restricted to the exterior at \(s=1\) gives
\eqref{eq:weighted-collar-exterior}. In particular \(V_j^\pm\)
can be negative.

At each path parameter a single-patch change alone is a smooth
metric on the same closed oriented \(M\). Its two total chiral
increments agree by signature, and agree with the corresponding
increments on that patch. The same reasoning applies after the
single patch is capped with its marked conformally flat disk.
Thus if
\(e+\sum_j\min_\pm U_j^\pm<2T\), the simultaneous joined path,
product stretching and compactification of
Theorem~\ref{thm:simultaneous} yield a positive class on the
original \(M\). Its energy is at most
\begin{equation}\label{eq:weighted-collar-gain}
 e-\sum_j\max_\pm\{\lambda^\pm(B_j)-V_j^\pm\}.
\end{equation}
In particular, the sign used to estimate the gain may vary by
disk. An actual single-patch exterior change \(o_j^\pm\)
satisfies \(o_j^+-o_j^-=\lambda^+(B_j)-\lambda^-(B_j)\)
by the same signature identity.

\paragraph{Finite measured selection.}
For each protected core suppose there are \(N_j\) candidate
annuli of width two, with the preceding nesting and disjointness.
Normalize each to the annulus of
Proposition~\ref{prop:disk-collar-comparison}. Define
\(W_j^2=\sum_\ell w_{j\ell}^2\),
\(L_j^2=\sum_\ell\tau_{j\ell}^2\),
\(P_j^2=\sum_\ell d_{j\ell}^2\), and \(a_j=\max_\ell a_{j\ell}\).
Then \eqref{eq:weighted-collar-allowance} and Cauchy--Schwarz give
\[
 \sum_\ell(w_{j\ell}^2+\beta_{j\ell}^2)
 \le \mathcal D_j
 =W_j^2+C\bigl[a_j^2(W_j^2+L_j^2)
                +P_j(W_j+L_j)+P_j^2\bigr].
\]
The physical curvature ceiling on the containing ball controls
\(W_j^2\). In a common parent chart the \(L_j^2\) term is the
integral of \(|DG|^2\) on the union of candidate annuli, by
critical scaling, and \(P_j^2\) is bounded by a fixed multiple
of the integral of \(|h|^2|x|^{-4}\) there. Neither of these
coordinate estimates follows from the physical curvature ceiling.

Choose the candidate minimizing \(w_{j\ell}^2+\beta_{j\ell}^2\)
on each core. For each sign and on either region the surcharge
is at most \(2\sqrt S\,\beta+\beta^2\le S+2\beta^2\), with
\(S\le w^2/4\). The summed joined and exterior allowances
are therefore each at most \(\Gamma=2\sum_j\mathcal D_j/N_j\).
All chosen disks contain the corresponding \(C_j\).
If \(e+\Gamma<2T\) and
\(\sum_j\max_\pm\lambda^\pm(C_j)>\Gamma+(e-I)\),
then \eqref{eq:weighted-collar-gain} yields a strict same-type
comparison. A known charge bound on the original disks transfers
to their smaller cores after subtracting the actual lost charge.
The cost and the protected charge are independent measurements on
the same finite family.
